\section{常见实验问题与适用范围}
\label{sec:failure-boundaries}

只有将失败或异常结果与组成、实验条件及产物表征一并记录，才能判断方法的适用范围。本节讨论相竞争、污染与挥发、气氛与晶型、资料缺失以及环境、健康与安全（EHS）限制；内容限于已发表条件，不将这些边界改写为新的实验方案。

\paragraph{相区与相竞争。}
Ba--Y硅酸盐系列的系统性助熔剂法生长比较表明，相形成结果必须结合组成和路线变量解读，不能孤立提取某一热处理条件作为通用条件\cite{wierzbickawieczorek2017high}。后续研究通过非化学计量结构的重新审定、陶瓷组成系列和多种衍射方法，揭示了调制结构、第二相以及竞争相或玻璃区\cite{yamane2024synthesis,yamane2024microstructure,gulay2024navigation}。这些结果仅界定各自考察的组成范围，不能证明未覆盖区域不存在其他相。机械化学研究仅提示预活化可能影响后续相形成，不能将某个负结果解释为目标相的固有相界\cite{tzvetkov2001effects}。因此，表~\ref{tab:synthesis-conditions}\CJKecglue{}和表~\ref{tab:synthesis-conditions-b}\CJKecglue{}中的产物、表征和合成条件必须出自同一文献，杂相或玻璃也不能脱离组成和热历史单独解释。

\paragraph{污染、挥发与助熔剂边界。}
无Zn的Ba--Y陶瓷研究将第二相与玛瑙来源的\(\mathrm{SiO_2}\)污染关联到具体处理流程，说明表观配比偏离应先检查原料和处理过程，不能自动视为真实固溶的证据\cite{yamane2024microstructure}。Y--Si--O体系的提拉法研究可提供熔体、坩埚、气氛和提拉条件的参照，但现有资料不足以判断目标相是否发生Zn或Si挥发\cite{shoudu1999czochralski}。含铅、含氟助熔体系曾用于Zn--Si--O晶体生长，但这些早期条件仅用于比较缓冷、分离和助熔剂作用，不能作为目标相的材料选择依据\cite{giess1982zn2sio4}。因此，应分别记录研磨介质带入、开放体系的物质损失以及助熔剂残留或相选择；文献未给出烧前后组成或残留分析时，本文不指定具体的失败机理。

\paragraph{气氛与晶型。}
BaZn$_2$Si$_2$O$_7$的低温单斜相与高温正交相由X射线和中子衍射区分，说明原位高温相与冷却后产物不可互换\cite{lin1999phase}。高温衍射还显示，Zn\slash Co取代以及Ba\slash Sr--Zn--Si\slash Ge固溶体组成与相稳定性或晶型边界联动；组成、温度和观察时点缺少任一项，均不足以定义稳定范围\cite{kerstan2013bazn2si2o7,thieme2016negative}。相关Ba--Zn--Si陶瓷的抗还原报道仅说明该化合物在已发表条件下经过气氛耐受性考察，不能外推为目标相的还原稳定性\cite{zou2019anti}；其他热膨胀与高温衍射比较也仅作为热致变化的参照\cite{kerstan2012thermal}。

\paragraph{资料缺失与EHS。}
部分文献仅给出合成方法或结构结论，未同时报道配比、温度--时间、气氛、坩埚和冷却等完整信息\cite{ababaikeri2024ba5y12zn,ababaikeri2024ba3zn4si4o15,kaiser2002crystal,cai2024optimized,zhao2026crystal}。这表示本文尚未取得相应资料，既不等于原文未作报道，也不等于实验中没有该步骤；表~\ref{tab:synthesis-conditions}\CJKecglue{}以“—”标记并保留这一边界。资料越不完整，可借鉴的具体条件越少。EHS边界同样仅按已发表条件描述：含Ba粉体、高温挥发性物质、早期文献中的Pb\slash F助熔剂条件以及坩埚相容性和废物流均需单独审核，其中Pb\slash F路线尤其不能由早期实例改写为推荐方案\cite{giess1982zn2sio4}。任何复用均须经过独立的EHS、容器相容性和废物分流审核；本文不补写具体操作参数，也不对未报道的通风、防护、设备和废物分类作合规推定。

上述四类边界不能合并为一个未经验证的失败机理。相竞争仅在已发表的组成和热历史范围内成立；污染或物质损失需要介质、烧前后组成或残留分析的支持；气氛和晶型判断必须保留观察温度及冷却后的状态；资料缺失仅说明当前认识的限度。因此，文献仅报道“伴随”或仅给出结构结果时，正文使用“伴随”或“一致”，而不用“导致”。这样既保留了负结果和成功记录，也避免将未报道项、工艺异常和真实相区混为一谈。
